\chapter{Итоговые проекты}\label{ch:proekty}

\textbf{Что вы сделаете:} соберёте всё вместе: от двух дронов до восьми. Задачи идут по нарастающей. Общий критерий для всех: \textbf{«Проверить» не показывает ни ошибок, ни замечаний} (сведения допускаются). Дроны летают в левой нижней четверти, клетки 0--5 по столбцам и строкам.

\begin{zametka}[Общие правила]
\begin{itemize}
  \item Стартовые клетки соседних дронов не ближе двух клеток (80\cm).
  \item Каждый дрон начинает с \kod{calibrate(5)} и \kod{takeoff(...)}, потом первая \kod{code} указывает его стартовую клетку.
  \item Расстояние между дронами не меньше 60\cm в любой момент. Если пути пересекаются, разводите дронов по высоте или по времени.
  \item Скорость перехода не больше 100\cm/с, высота не больше 150\cm.
  \item Заканчивайте программу посадкой.
\end{itemize}
\end{zametka}

\section{Проект 1. За лидером (2 дрона)}

Дрон 1 (80\cm) облетает квадрат 13, 16, 52, 49 и возвращается в клетку 13. Дрон 2 (150\cm) стартует из клетки 37 и повторяет тот же маршрут на один переход (4 секунды) позже, а в конце садится в клетку 37.

\emph{Подсказка.} Дрон 2 сразу выходит в клетку 13 за 4 секунды: пока лидер идёт в клетку 16, он занимает его место. Тогда он всё время на один шаг позади.

\section{Проект 2. Лесенка (3 дрона)}

Три дрона в клетках 13, 15 и 17 на разных высотах сдвигаются по кругу: 1 на место 2, 2 на место 3, 3 на место 1.

\emph{Подсказка.} Дрон, летящий в самый дальний угол, пройдёт над двумя другими. Подберите три высоты от 20 до 150\cm с зазором не меньше 60\cm.

\section{Проект 3. Четыре угла (4 дрона)}

Четыре дрона стоят по углам квадрата 5$\times$5 клеток (клетки 0, 4, 52, 48) и обходят его по кругу: каждый переходит в следующий угол и делает полный оборот.

\emph{Подсказка.} Все дроны движутся одновременно в одну сторону, на разные высоты разводить не нужно. Сторона 160\cm, хватит 4 секунд на переход. Вариант с циклом лежит в файле \texttt{t10\_loop.py}.

\section{Проект 4. Пятёрка (5 дронов)}

Четыре дрона по углам квадрата идут по кругу на два угла вперёд, а пятый в центре (клетка 26) поднимается на 120\cm и опускается на 60\cm.

\emph{Подсказка.} Центр отстоит от угловых на 113\cm, а в середине стороны угловые проходят в 80\cm от него: этого хватает без разведения по высоте.

\section{Проект 5. Ряды меняются (6 дронов)}

Шесть дронов стоят в двух рядах по три: нижний ряд в клетках 13, 15, 17, верхний в клетках 49, 51, 53. Ряды меняются местами. Нижний ряд летит на 80\cm, верхний на 150\cm.

\emph{Подсказка.} Дроны одного столбца летят навстречу друг другу, поэтому им нужны разные высоты. Соседние столбцы разделены 80\cm.

\section{Проект 6. Семёрка (7 дронов)}

Шесть дронов стоят шестиугольником вокруг клетки 26 (клетки 2, 16, 40, 50, 36, 12), седьмой в центре. Шестёрка по кольцу сдвигается на два места, а центральный мигает белым пять раз.

\emph{Подсказка.} Соседние дроны кольца на высотах 90 и 130\cm, чтобы разойтись при сдвиге. Центральный висит ниже: 60\cm.

\section{Проект 7. Хоровод (8 дронов)}

Восемь дронов стоят по кольцу в клетках 0, 2, 4, 28, 52, 50, 48, 24. Дроны с чётными номерами (на 150\cm) идут по кольцу в одну сторону, с нечётными (на 80\cm) навстречу. Каждый проходит два места и садится в центр клетки.

\emph{Подсказка.} Соседние дроны на разных высотах, поэтому встречные разминутся. Кольцо задайте списком и берите клетки по индексу с остатком от деления на 8.

\otvety{16}
